\section{The correction-factor conjecture in an eventual dense square regime}\label{sqt:sec:conjecture}
Canfield--McKay's Conjecture~1 \cite[Eq.~(1.4)]{CM} defines a correction $\Delta(m,s;n,t)$ by
\begin{align}\label{sqt:eq:CMdelta}
M(m,s;n,t)={}&G_{\rm Good}(m,s;n,t)
 \left(1+\frac1m\right)^{(m-1)/2}
 \left(1+\frac1n\right)^{(n-1)/2}\notag\\
&\times\exp\left(-\frac12+\frac{\Delta(m,s;n,t)}{m+n}\right),
\end{align}
where $ms=nt$, $M$ is their table count, and
\[
G_{\rm Good}(m,s;n,t)=
\frac{\binom{n+s-1}{s}^{m}\binom{m+t-1}{t}^{n}}
 {\binom{mn+ms-1}{ms}}.
\]
They conjecture $0<\Delta<2$ for all positive integer admissible parameters. This $G_{\rm Good}$ differs from the Gaussian prefactor $G_n$ of \eqref{sqt:eq:leading}.

\begin{corollary}\label{sqt:cor:delta}
On each fixed compact interval $I\subset(0,\infty)$, uniformly over square tables with integer line sum $s$ and $\lambda=s/n\in I$,
\begin{equation}\label{sqt:eq:deltaasymp}
\Delta(n,s;n,s)=\frac1n\left(\frac{67}{6}+\frac5{3u}\right)+O_I(n^{-2}),
\qquad u=\lambda(1+\lambda).
\end{equation}
Consequently $0<\Delta(n,s;n,s)<2$ for all sufficiently large $n$, uniformly over those margins. In particular,
\begin{equation}\label{sqt:eq:deltadensityone}
\Delta(n,n;n,n)=\frac{12}{n}+O(n^{-2}).
\end{equation}
\end{corollary}
\begin{proof}
Set $\mathcal L(n,\lambda)=\mathcal G(n,\lambda)e^{B_0(u)}$. The standard Stirling expansion of the three factorial ratios in $G_{\rm Good}$ gives, uniformly on $I$,
\begin{equation}\label{sqt:eq:goodstirling}
\log\frac{G_{\rm Good}(n,s;n,s)e^{1/2}}{\mathcal L(n,\lambda)}
 =\frac{g_2(u)}{n^2}+O_I(n^{-4}),
\end{equation}
where
\begin{equation}\label{sqt:eq:g2}
g_2(u)=\frac4{45}+\frac1{12u}+\frac1{60u^2}+\frac1{180u^3}.
\end{equation}
Here is an explicit way to see the coefficient. If
$B(N,\lambda)=\binom{(1+\lambda)N-1}{\lambda N}$, write
\[
\log B(N,\lambda)=N\log H(\lambda)-\tfrac12\log(2\pi N u)
 +\frac{b_1(u)}N+\frac{b_3(u)}{N^3}+O_I(N^{-5}),
\]
with
\[
b_1(u)=-\frac{1+1/u}{12},\qquad
b_3(u)=\frac{1+3/u^2+1/u^3}{360}.
\]
The constant in this display needs the factor $1/(1+\lambda)$ in
$\binom{(1+\lambda)N-1}{\lambda N}=\binom{(1+\lambda)N}{\lambda N}/(1+\lambda)$; it is already included in $-\tfrac12\log(2\pi N u)$. Since
$\log G_{\rm Good}=2n\log B(n,\lambda)-\log B(n^2,\lambda)$,
its $n^{-2}$ coefficient is $2b_3-b_1=g_2$. Its constant is $2b_1=B_0-1/2$. This proves \eqref{sqt:eq:goodstirling}.

Taking logarithms in \eqref{sqt:eq:CMdelta} in the square case,
\[
\frac{\Delta}{2n}=\log\frac{T(n,s)}{G_{\rm Good}}-(n-1)\log(1+1/n)+\frac12.
\]
Use Theorem~\ref{sqt:thm:density}, \eqref{sqt:eq:goodstirling}, and
\[
(n-1)\log(1+1/n)=1-\frac{3}{2n}+\frac5{6n^2}+O(n^{-3}).
\]
The constants and $n^{-1}$ terms cancel. The remaining coefficient is
\[
P_2(1/u)-g_2(u)-\frac56=\frac{67}{12}+\frac5{6u}.
\]
Multiplication by $2n$ proves \eqref{sqt:eq:deltaasymp}. The leading coefficient is positive and bounded above and below on $I$, so its uniform remainder proves the eventual strict inequalities. Substituting $u=2$ gives \eqref{sqt:eq:deltadensityone}.
\end{proof}

This verifies the conjectured inequalities in one regime only: square tables, density in a fixed compact subset of $(0,\infty)$, and $n$ beyond a threshold that is not made explicit (Remark~\ref{sqt:rem:cmregime}). It does not settle all finite sizes, nonsquare tables, or densities tending to zero or infinity. Prior progress includes Greenhill--McKay \cite[Corollary~4.2]{GM}, which proves
$\Delta(m,s;n,t)\sim5(s+t)/(6st)$ in the sparse regime
$st=o((mn)^{1/5})$. Their preceding Corollary~4.1 also contains $-3/(4n)-3/(4m)$, hence $-3/(2n)$ for squares, under sparse assumptions. Those assumptions exclude fixed positive square density. Thus the numerical first coefficient has a sparse-regime antecedent; the present assertion concerns its analytically controlled dense regime and the next correction.

The explicit prefactor in the original conjecture also suggests a term $-3/(2n)$ on expansion. It did not prove that coefficient, because the unknown term $\Delta/(2n)$ could have contributed at the same order. The present corollary shows $\Delta\to0$ in the stated regime.

\begin{remark}[{[write]} The conjecture and the claim, quoted; what is proved, 6~October 2026]\label{sqt:rem:cmregime}
(a) \emph{The conjecture.} Conjecture~1 of \cite{CM} (arXiv:math/0703600v2,
p.~4) reads: ``Consider a 4-tuple of positive integers $m,s,n,t$ such that
$ms=nt$. Define $\Delta(m,s;n,t)$ by'' \eqref{sqt:eq:CMdelta}. ``Then
$0<\Delta(m,s;n,t)<2$.'' Canfield and McKay record it as proved ``(a) for
$m=n\le9$, using the exact values'' of the Ehrhart polynomials; ``(b) for
sufficiently large $m,n$ when $st=o\bigl((mn)^{1/5}\bigr)$'', using
Greenhill--McKay's asymptotics \cite{GM}; and (c) for several thousand
tuples with $m,n\le30$.

(b) \emph{The claim, as delivered.} The abstract said the density extension
``proves an eventual square case of the Canfield--McKay correction-factor
conjecture''; this section was titled ``An eventual square case of a
correction factor conjecture''; and the sentence after
Corollary~\ref{sqt:cor:delta} began ``This is an eventual
compact-positive-density \emph{square} case of the conjecture, with no
explicit starting index.'' The write narrowed all three on 6~October 2026
(this is the first wording).

(c) \emph{What is proved.} Corollary~\ref{sqt:cor:delta} gives: for every
compact $I\subset(0,\infty)$ there is $n_0(I)$ such that
$0<\Delta(n,s;n,s)<2$ for all $n\ge n_0(I)$ and all integers $s$ with
$s/n\in I$. The threshold $n_0(I)$ comes from the implied constants of
Theorem~\ref{sqt:thm:density} and of the imported theorems and is not
effective, so the corollary verifies the conjecture at no explicitly named
tuple, and it says nothing when $m\ne n$. For each fixed rational density
$\lambda$ it covers the tuples $(n,\lambda n;n,\lambda n)$ with $n$ large,
but not uniformly as $\lambda\to0$ or $\lambda\to\infty$. Square tuples
covered neither by the corollary nor by Canfield and McKay's cases~(a)
and~(b) (their case~(c) is a finite set with $m,n\le30$ that they do not
list) therefore remain for every choice of the $I$'s: sizes $10\le
n<n_0(I)$ within the regime, and, for large $n$, densities $s/n\to0$ with $s$
not $o(n^{1/5})$ (at $m=n$, $s=t$, case~(b) needs $s^2=o(n^{2/5})$), and
$s/n\to\infty$. The corollary settles one regime of the conjecture: not the
conjecture, and not its square case.

(d) \emph{Exact counts.} From the OEIS values (Remark~\ref{sqt:rem:oeis}) and
the definition~\eqref{sqt:eq:CMdelta}, evaluated in 60-digit floating
arithmetic (not interval arithmetic):
\begin{center}\footnotesize\setlength{\tabcolsep}{3pt}
\begin{tabular}{@{}r*{13}{c}@{}}
\toprule
$n$ & 1&2&3&4&5&6&7&8&9&10&11&12&13\\
\midrule
$\Delta(n,n;n,n)$ & 1&1.416&1.291&1.070&0.950&0.869&0.805&0.751&0.704&0.663&0.627&0.594&0.565\\
$n\Delta$ & 1&2.83&3.87&4.28&4.75&5.21&5.63&6.01&6.33&6.63&6.89&7.13&7.35\\
\bottomrule
\end{tabular}
\end{center}
($\Delta(1,1;1,1)=1$ exactly.) Every value lies in $(0,2)$, at distance at
least $0.56$ from both ends; the cases $n\le9$ are within Canfield and
McKay's case~(a). Corollary~\ref{sqt:cor:delta} predicts $n\Delta\to12$.
The products $n\Delta$ increase through $n=13$ and are still far from $12$;
the ratio $\Delta/(12/n)$ rises from $0.36$ at $n=4$ to $0.61$ at $n=13$.
Richardson extrapolation of order $k$ (the model
$n\Delta=c_0+c_1/n+\cdots+c_k/n^k$ fitted exactly at $n=13-k,\ldots,13$)
gives $c_0\approx9.94$, $10.98$, $11.43$, $11.66$, $11.78$ for $k=1,\ldots,5$.
This is consistent with the limit $12$, and it is not a proof of anything;
in particular it bounds no $n_0(I)$.

(e) \emph{The sparse end, formally.} As $\lambda\to0$ the coefficient
in~\eqref{sqt:eq:deltaasymp} behaves like $5/(3un)\sim5/(3\lambda n)=5/(3s)$,
which is the square case $s=t$ of the sparse asymptotic
$\Delta\sim5(s+t)/(6st)$ attributed above to Greenhill--McKay. The two
regimes are thus formally compatible; substituting $\lambda\to0$ into a
fixed-density formula proves nothing about the transition (the source's
Question~4 in Section~\ref{sqt:sec:outlook}; Question~\ref{sqt:q:regime}).
\end{remark}
\begin{writenote}[Independent check of the write, 6~October 2026]
An adversarial check made by the intake after the write (commit
\file{d63de7211}) re-derived the statements marked [write] with its own
code, after fetching again the OEIS entry A110058 and its b-file (revision
and data as quoted), Canfield--McKay's arXiv:math/0703600v2 and Isaev's
arXiv:2508.16952v2. \emph{Lemma~\ref{sqt:lem:sequential}.} The subsequence
argument was re-read and holds; against the rendered preprint, the region~(3.1)
(bounds $(1+\lambda)^{-1}(mn)^{-1/2+2\varepsilon}$,
$(1+\lambda)^{-1}n^{-1/2+\varepsilon}$, $(1+\lambda)^{-1}m^{-1/2+\varepsilon}$
on open half circles), Theorem~3 and the convention of p.~4 are as the
lemma uses them, and~(4.1) at $m=n$ is exactly $I_{0,\lambda}$; the left side
of~(1.1) at $m=n$ is $\tfrac23(4+1/u)$, equal to $3$ at $\lambda=1$ (SymPy).
The paragraph after the lemma matches the two places where the minor-arc
integral enters the proof of Lemma~\ref{sqt:lem:transfer}.
\emph{Remark~\ref{sqt:rem:uniformity}.} Every quotation and every bound
(the factors $\lambda(1+\lambda)^{-1}>(\log n)^{-1}$, $A=O(\log n)$,
$N=\lceil6000(1+\lambda)\rceil$, $\delta=2\pi/N$, the three pieces with
$O(e^{-n^{1-\varepsilon}})I_0$ and $e^{-n^\varepsilon}I_0$, and the omitted
proofs of Lemmas~3 and~4) was found on pp.~15--20; Isaev's Theorem~1.2 and
the inequality $\overline\Delta_V(f,X)\le\Delta_V(f)$ (``Clearly'') were
found as Section~\ref{sqt:sec:provenance} says.
\emph{Remark~\ref{sqt:rem:cmregime}.} The quotations of Conjecture~1 and of
the record of proved cases are verbatim, and the three first wordings are
those of the delivered files. From the fresh b-file, at $80$ digits, all
thirteen values of $\Delta(n,n;n,n)$ and $n\Delta$ in the table, the margin
($\min=0.5651$ at $n=13$), the ratios $0.36$ and $0.61$, both
monotonicities of Question~\ref{sqt:q:monotone} and the Richardson values
$9.94$, $10.98$, $11.43$, $11.66$, $11.78$ were reproduced;
$P_2(1/u)-g_2(u)-\tfrac56=\tfrac{67}{12}+\tfrac5{6u}$ and the expansion of
$(n-1)\log(1+1/n)$ were rechecked in SymPy.
\emph{Remark~\ref{sqt:rem:oeis}.} Revision~\#30, its name, data, link and
Pak's formula are as quoted, Pak's main terms are exactly
$\log L_n-\tfrac12\log(4\pi)-\tfrac14$ (SymPy), and
$\tfrac12\log(4\pi)+\tfrac14=1.515512123\ldots$.
\emph{Remark~\ref{sqt:rem:transseries}.} Parts~(1) and~(3) hold (the master
equation is exact and the triangular recurrence gives $d,e,f,g$). One
statement was too strong: (2) said that the leading balance is not of the
monomial--logarithmic type of \file{plt:thm:lw-template}, but after
$G=\sqrt{F_2/a}$ the reversion is a formal instance of that theorem, as for
the sibling reports \file{a089479-fixed-permanent-matrices} and
\file{a222959-zero-slope-matrices}; corrected with a dated note keeping the
first wording, and the instance added to~(3). \emph{Provenance.} The archive
(669{,}911 bytes, the SHA-256 quoted, 35 files, 34 manifest entries, 980
delivered lines, 24 pages) and the counts of 97 labels and 87 references
were confirmed. For the record, the intake's Richardson values $-1.47$ and
$6.5$ for $D_1$, $D_2$ match fits of order two or three at the top of the
range; the exact fit through all eight values $6\le n\le13$ gives
$-1.4964$ and $6.9139$. No other defect was found.
\end{writenote}
